\section{Method}

\subsection{Dense \Htwo-style residual baseline}

The \Htwo-style baseline fits one dense residual model over all transition outputs:
\[
    \hat \Delta_{\mathrm{dense}}(s,a) = \phi_{\mathrm{dense}}(s,a)^\top B,
\]
where $\phi_{\mathrm{dense}}$ includes normalized state/action features, squares, action interactions, and selected pairwise products.
At evaluation time the corrected prediction is
\[
    \hat \y_{\Htwomath}(s,a) = f_0(s,a) + \hat \Delta_{\mathrm{dense}}(s,a).
\]
The baseline is intentionally strong: it can use many cross terms and is fit on all source-city transitions.
It is nevertheless an \Htwo-style dense residual transfer surrogate rather than a full reproduction of the original H2O+ policy-learning stack; \Cref{tab:h2oplus-fidelity} documents the implementation differences.

\subsection{Causal-factored residual transfer}

\ours decomposes the residual by output mechanism:
\[
    \hat \Delta_j(s,a) = \phi_j(s,a)^\top \beta_j, \qquad j \in \{1,\ldots,9\}.
\]
The feature map $\phi_j$ is chosen according to the output mechanism.
Headway and gap residuals use current headways, target headway, speed, action, and time.
Waiting and dwell residuals use waiting passengers, target headway, speed, station fraction, action, and time.
Speed residuals use current speed, event time, and station position.
Reward residuals use policy-relevant headway deviations, waiting passengers, gap, speed, action, and time.
The corrected prediction is
\[
    \hat \y_{\oursmath}(s,a)
    = f_0(s,a) +
      \left(\hat\Delta_1(s,a),\ldots,\hat\Delta_9(s,a)\right).
\]
All reported models use ridge-regularized sufficient statistics, so the experiment can be reproduced without stochastic neural training.

\subsection{Theoretical motivation}

The theory we need is not a proof of causal identification.
Instead, it explains why a mechanism-factored residual and a guarded source ensemble are natural for cross-city transfer when target transition labels are unavailable.
Let $x=(s,a)$ and let
\[
    \delta_c(x)=f_c(x)-f_0(x)
\]
denote the simulator residual for city $c$.
For a residual predictor $g$, define the target risk
\[
    R_t(g)=\E_{x\sim P_t}\left[\|\delta_t(x)-g(x)\|_2^2\right],
\]
where $P_t$ is the target-city decision-state distribution induced by the offline validation set.
Partition the nine outputs into mechanism groups $\mathcal{M}_1,\ldots,\mathcal{M}_K$, such as headway/gap, demand/dwell/reward, and speed.
For $g=(g_1,\ldots,g_K)$ and $\delta_c=(\delta_{c,1},\ldots,\delta_{c,K})$, squared-error risk decomposes exactly:
\[
    R_t(g)=\sum_{k=1}^{K}
    \E_{x\sim P_t}\left[
        \|\delta_{t,k}(x)-g_k(x)\|_2^2
    \right].
\]
Thus each mechanism contributes an auditable term to total error.
This does not imply that the groups are causally identified, but it gives a formal reason to fit, validate, and report mechanism-wise residual errors rather than only a dense total error.

\paragraph{Parent-set invariance assumption.}
The operative assumption is weaker than full causal identification.
For each output mechanism $k$, there exists a small parent set $Pa_k(x)$ such that the residual map $\delta_{c,k}$ varies across cities mainly through mechanism-specific coefficients, context bias, or source weights, while irrelevant city-specific nuisance features do not improve target risk after source validation.
The sparse parent sets in \ours encode this assumption for headway/gap, demand/dwell/reward, and speed outputs.
The assumption can fail: if the target generator or real operation shifts an omitted parent, the factored model can underfit.
This is why the experiments include wrong-grouping, action/time-only, dense matched-feature, source-subset, and generator-bias diagnostics, and why the paper avoids claiming causal identification from observational static data alone.

\paragraph{Mechanism-shift bound.}
Let $\bar\delta_{w,k}=\sum_{c\in\mathcal{C}_{\mathrm{src}}} w_c \delta_{c,k}$ be the weighted source residual for mechanism $k$.
Suppose the fitted mechanism predictor satisfies
\[
    \left\|\hat g_k-\bar\delta_{w,k}\right\|_{L_2(P_t)}\le \epsilon_k,
    \qquad
    \left\|\bar\delta_{w,k}-\delta_{t,k}\right\|_{L_2(P_t)}\le \Delta_k(w),
\]
where $\epsilon_k$ is estimation/approximation error and $\Delta_k(w)$ is the cross-city mechanism shift left after source weighting.
By the triangle inequality,
\[
    \sqrt{R_t(\hat g)}
    \le
    \left(
        \sum_{k=1}^{K}
        \left[\epsilon_k+\Delta_k(w)\right]^2
    \right)^{1/2}.
\]
The bound clarifies the intended advantage: if only some mechanisms shift strongly across cities, a factored model can localize transfer error to those mechanisms.
A dense residual can still perform well, but its error is harder to attribute and may mix source-specific headway, demand, dwell, and speed shifts in a single parameter block.

\paragraph{Source-ensemble risk bound.}
For source-city experts $g_c$, the ensemble predictor $g_w=\sum_c w_c g_c$ has a convexity bound under squared loss:
\[
    R_t(g_w)
    =
    \E_{P_t}\left[
        \left\|
        \sum_c w_c(\delta_t-g_c)
        \right\|_2^2
    \right]
    \le
    \sum_c w_c R_t(g_c),
    \qquad
    w_c\ge0,\ \sum_c w_c=1.
\]
The unavailable quantity is $R_t(g_c)$.
The target-label-free source-weighting rule therefore uses unlabeled target summaries as a proxy for which source-city expert is likely to have low target risk.
The dominance guard $\max_c w_c(t)\ge0.5$ is a conservative negative-transfer control: an ensemble is allowed to override the rigid factored residual only when the unlabeled target summary identifies at least one clearly close source; otherwise the method falls back to the rigid source-pooled model.

\paragraph{Source-only selector bound.}
Let $q$ index a candidate target-summary feature family, and let $S(q)$ be its leave-one-source validation score.
Let $R_t(q)$ denote the target risk of the full guarded procedure using feature family $q$.
If the source-fold score approximates the target ranking up to
\[
    |R_t(q)-S(q)|\le \gamma + \varepsilon
    \quad\text{for all candidate families }q,
\]
where $\gamma$ is the source-to-target representativeness gap and $\varepsilon$ is finite-sample validation error, then the selected family
\[
    \hat q=\arg\min_q S(q)
\]
satisfies
\[
    R_t(\hat q)
    \le
    \min_q R_t(q)+2(\gamma+\varepsilon).
\]
The condition is not guaranteed by the data; it is the assumption being tested by source-subset and feature-family diagnostics.
It is nevertheless important methodologically: feature-family selection can be done without target transition labels, and any failure mode is visible as a source-representativeness problem rather than target-label tuning.

\subsection{Unlabeled target source selection}

For each city we compute a static descriptor $z_c$ from speed quantiles, headway quantiles, route-distance quantiles, peak-demand hour, and the 24-hour demand-share vector.
For a target city $t$ and source city $c$, the distance is computed after feature-wise standardization over the source-plus-target set:
\[
    d(c,t) =
    \frac{1}{\sqrt{m}}\left\|
    \frac{z_c-\mu}{\sigma} -
    \frac{z_t-\mu}{\sigma}
    \right\|_2.
\]
The source weight is
\[
    w_c =
    \frac{\exp(-d(c,t)/T) + \lambda \max_{c'} \exp(-d(c',t)/T)}
         {\sum_{k \in \mathcal{C}_{\mathrm{src}}}
          \left[\exp(-d(k,t)/T) + \lambda \max_{c'} \exp(-d(c',t)/T)\right]}.
\]
The default values are $T=1.0$ and $\lambda=0.05$.
The source sufficient statistics are rescaled so that these weights control the effective number of source transitions rather than simply reflecting raw city size.

\subsection{Rigid-first source-city ensemble}

The full-module variant starts from the causal-factored ridge prediction rather than from the uncorrected simulator:
\[
    \hat \y_{\mathrm{rigid}}(s,a)
    = f_0(s,a) + \hat \Delta_{\oursmath}(s,a).
\]
It then fits a second residual on the remaining source error,
\[
    e_c(s,a)=\y_c^{\mathrm{target}}-\hat \y_{\mathrm{rigid}}(s,a),
\]
using dense features, deterministic local-graph features, or their concatenation.
A candidate correction has the form
\[
    \hat \y_k(s,a)
    =
    \hat \y_{\mathrm{rigid}}(s,a)
    + \alpha \eta_k(s,a),
    \qquad
    \alpha \in \{0.25,0.50,0.75,1.00\}.
\]
For the source-city ensemble candidate, the second residual is fit separately on each source city and mixed by target similarity:
\[
    \eta^{\mathrm{ens}}_k(s,a;t)
    =
    \sum_{c\in\mathcal{C}_{\mathrm{src}}} w_c(t)\eta_{k,c}(s,a).
\]
The source-summary feature set is not chosen by target performance.
For each candidate summary family (simulator-output only, observation only, observation plus simulator output, local-only, and full summary), the same leave-one-source folds below are used to score the selected residual candidate.
The summary family with the best source-fold score is then used to compute $w_c(t)$ for the held-out target.
None of these variants uses target transition labels.

\paragraph{Source-summary feature-set selector.}
The guarded ensemble uses the following target-label-free selection procedure.
\begin{enumerate}[leftmargin=1.5em]
    \item Define the candidate summary families listed in \Cref{tab:target-summary-families}.
    \item For each candidate family, run the same leave-one-source folds used for residual-candidate selection.
    \item Within each fold, select the residual candidate by source-fold total-MSE ratio to rigid \ours, subject to the ensemble dominance guard.
    \item Score the summary family by the selected candidate's mean source-fold ratio; lower is better.
    \item Break ties by preferring simulator-output summaries, then observation-plus-simulator, observation-only, full, local-only, and uniform weighting.
    \item Freeze the selected summary family before target evaluation and compute target source weights from unlabeled target summaries only.
\end{enumerate}

Candidate selection uses leave-one-source validation.
For each held-out source city $h$, we train on $\mathcal{C}_{\mathrm{src}}\setminus\{h\}$, evaluate on $h$, and record each candidate's total-MSE ratio to rigid \ours.
The selected candidate must have stable source-fold gains.
Ensemble candidates are additionally guarded: they may override the conservative rigid/global candidate only when the target similarity weights have a dominant source, $\max_c w_c(t)\ge 0.5$.
This guard prevents a diffuse source mixture from overriding rigid \ours when the unlabeled target summary does not identify a close source city.

\subsection{Traffic-signal phase-transfer stress test}

For the RESCO stress test, the same transfer principle is applied to phase-local traffic-signal models.
For each traffic light and candidate green phase, a generic simulator prior predicts the next local queue cost.
The dense baseline fits one pooled residual over all phase-local features, while the CFCMT variant decomposes the residual into queue-propagation, downstream-spillback, and phase-pressure feature groups.
At each decision time, model-based policies evaluate the feasible SUMO-defined green phases and select
\[
    a^\star_t
    =
    \arg\min_{a\in\mathcal{A}(i)}
    \hat c(x_{i,t},a),
\]
where $\hat c$ is the corrected next-interval local queue cost.

Classical pressure control is a hard baseline in this setting.
We therefore report both a phase-pressure rule and a MaxPressure-style rule that chooses the feasible phase with the largest upstream-minus-downstream queue relief.
The reported guarded CFCMT controller treats MaxPressure as a safety prior.
It first computes the CFCMT phase-MPC action and the MaxPressure action, evaluates both under the CFCMT cost model, and permits the CFCMT action only when its predicted cost is lower by a fixed margin.
The guard also falls back to MaxPressure for pre-specified target-static regimes where source transfer was unstable in source-only diagnostics: large regular networks, very high-demand single-intersection cases, and sparse low-demand targets.
These guard inputs are target descriptors and current queues; they do not use target closed-loop rollout labels.
This is the control analogue of the bus-holding conclusion: a transferred mechanism model is used to instantiate and gate a physically meaningful prior, not to replace that prior with an unconstrained residual action rule.

\subsection{One-step policy proxy}

For policy-level validation, each method evaluates a discrete hold set
$\{0,15,30,45,60\}$ seconds under a set of headway-stress states.
The policy selects the action with the highest predicted reward under its corrected transition model.
The selected action is then scored against the static-derived target transition.
We report reward, regret to an oracle action, mean absolute headway error, bunching rate, large-gap rate, and mean hold seconds.
This is a policy-level one-step lookahead benchmark, not a live SUMO rollout.
